% Methods: the hindcast audit of the relief grids (H1; A7). Design: D-030, D-032, D-033,
% D-050, D-067, D-070, D-072; PREREG section 9 and Amendments 6, 7, 10 and 11.
% Every number comes from \metric{...}.

\subsection*{Hindcast audit}
\emph{Grid and extents.} The audit runs on the \metric{params.stage2_res_m}~m water-frequency grid, extended to the domain with a margin of \metric{params.hindcast_margin_deg}\textdegree{}. Each relief grid was interpolated bilinearly between its own cell centres in its native longitude--latitude grid, and a cell next to a missing native value has no grid state; nearest-neighbour resampling, which gives every grid a staircase shoreline one native cell wide, is a sensitivity. For each epoch and level bin, a cell is valid with at least \metric{params.min_valid_obs} valid acquisitions and observed water with a frequency of at least \metric{params.contour_level}. As for the predictions, only water 8-connected through known water to the open-sea seed counts; unconnected water bordered only by known land is land, and unconnected water touching an unknown cell is unknown (pre-registration Amendment~7).

\emph{Known states outside the observed area} (pre-registration Amendments~6 and~10). Inside the domain, a connected group of cells outside the stage-2 cover is water if it holds an open-sea seed, land if it holds a land seed, and unknown otherwise; a connection between the two kinds of seed stops the analysis. The open-sea seed is the open-sea mask of the active zone outside the cover. Beyond the domain and in a one-cell rim inside it, a cell is land where it lies more than \metric{params.zone_dilate_m}~m from the \metric{params.highstand_label} near-maximum water of the active zone, mapped from Landsat with the Global Surface Water record where clear observations are too few\citep{pekel2016}, and unknown otherwise; unobserved cells count as possible water, and the state is never water, which would add agreement every grid gets trivially. This assumes that no evaluated level floods land beyond the observed near-maximum extent, the highest of the levelled record, and uses no relief grid. The land seeds are its land cells within \metric{params.beyond_inner_km}~km inside the domain edge. A cell outside the domain belongs to the sector of the nearest cell inside it.

\emph{Predicted extents and variants.} A grid's predicted extent at a level is its flood fill from the open-sea seed: the cells below the level that are 8-connected to the seed through such cells. All levels were computed at once from each cell's flood level, the minimax height of the paths joining it to the seed, so that a cell is water at level $L$ exactly when its flood level is below $L$. The total variant uses the grid on the DAHITI reference (Vertical datums). The shape variant subtracts the constant $c$ that minimises the summed absolute area bias over all epochs and levels, over known cells inside the domain in training blocks, searched in steps of \metric{params.shape_step_m}~m up to \metric{params.shape_max_m}~m either way (ties to the smallest $|c|$, then to the negative value), and refitted without each epoch for the leave-one-epoch-out check.

\emph{Metrics.} For each grid, variant, sector and level we computed the absolute and relative area bias, the intersection over union and the F1 score, counting errors at known cells outside the valid area but not correct predictions there, which would make agreement trivially high; values within the valid area are also reported. A waterline cell is known water with a known-land 4-neighbour. The shoreline displacement is symmetric: distances from each predicted waterline cell to the nearest observed one are pooled with those from each observed waterline cell to the nearest predicted one, over cells whose state on the other extent is known, and unmeasured cells are counted. We report the median and \metric{params.err_pct_hi}th percentile of the absolute distances and the signed median, positive where a grid predicts too much water; a grid without a waterline has infinite displacement. Per sector, waterline cells in the sector are measured against the nearest waterline anywhere. The noise floor is the same displacement between the extents of the two acquisition halves of each epoch and bin, each with its own valid area, and, at levels crossed twice, between E1 and E2 and E3 together.

\emph{Decision rules.} The shape variant decides H1, so that a datum mismatch alone cannot make a grid fail; the total variant is reported beside it, and a grid that fails only through its offset is named. A grid mispredicts in a sector if, at more than half of the sector's evaluable observed levels, its median absolute displacement there (epochs pooled) exceeds the \metric{params.err_pct_hi}th percentile of the noise-floor displacements at the same level. An observed level has an observed waterline in the sector, and is evaluable if at least one grid distance and one noise-floor distance were measured; a floor pooled over all levels is a sensitivity. Whether errors are largest in deltaic and vegetated sectors was tested on blocks. A block's value is the median absolute displacement of a grid's waterline cells in it, over all epochs and levels, for blocks with at least \metric{params.h1_min_block_n} distances. The statistic $T$ is the difference in median block value between deltaic or vegetated sectors and the rest, with a one-sided $p = (1 + \#\{T_{\mathrm{perm}} \ge T\})/(N + 1)$ from $N = \metric{params.permutations_n}$ permutations of the labels over blocks, at $\alpha = \metric{params.alpha}$. The grids share sources, so they are reported as partly dependent, not as independent replications. Seabed change between epochs, measured between E1 and later contours at levels crossed twice, is tested as an alternative explanation before any error is attributed to the grids.
